\documentclass[UTF8, 12pt, fontset=fandol, oneside]{ctexart}
\usepackage{researchnotes}

% 保留分册独立编译时的章号前缀。
\renewcommand{\thesection}{17.\arabic{section}}

\title{第十七章\quad 含参变量积分}
\author{}
\date{}

\begin{document}

\maketitle

含参变量积分在许多自然科学问题中常常碰到. 例如, 一个不稳定流速场可以表示成 $F(x,y,z,t),(x,y,z)\in D,t\in(a,b)$, 其中 $D$ 是 $\mathbb R^3$ 中的区域, $t$ 是时间变量. 一般说来, $F(x,y,z,t)$ 是 $(x,y,z,t)$ 的连续函数. 现设一个光滑双侧曲面 $S$ 位于 $D$ 内, 取定 $S$ 的一侧, 则在任意时刻 $t$, $F$ 在 $S$ 上的第二型曲面积分与流量密切相关. 对固定的时刻 $t$, 我们可以在 $S$ 上对 $F$ 进行第二型曲面积分. 当 $t$ 变化时, 该积分便是 $t$ 的函数. 一个自然的问题是: 该积分是否是 $t$ 的连续函数? 如果该积分是 $t$ 的连续函数, 则通过对 $t$ 积分就可计算出某时间段内流体通过 $S$ 的流量. 以后我们还要研究: 当 $F$ 关于 $t$ 可求偏导数时, 该积分关于 $t$ 的导数是否存在? 这些问题的解决就要用到含参变量积分的理论。

在本章中, 我们将重点讨论含有参数的函数的积分问题, 讨论的积分以定积分与广义积分为主. 读者可以将这些积分推广到其他积分 (如重积分、曲线积分和曲面积分等) 的情形。

\section{含参变量定积分}

设函数 $f(x,y)$ 在平面区域 $D=[a,b]\times[c,d]$ 上有定义. 若对于 $\forall x\in[a,b]$, 定积分 $\int_c^d f(x,y)\mathrm dy$ 存在, 则由此定义了区间 $[a,b]$ 上的函数
\[
  I(x)=\int_c^d f(x,y)\mathrm dy.
\]
我们称如此定义的函数为\textbf{含参变量定积分} (简称\textbf{含参变量积分}), 其中 $x$ 为参变量. 同理, 若对于 $\forall y\in[c,d]$, $J(y)=\int_a^b f(x,y)\mathrm dx$ 存在, 则也称 $J(y)$ 为含参变量定积分, 其中 $y$ 为参变量。

由于定积分是一个极限过程, 因此对含参变量积分的研究有点类似于对函数序列或函数项级数的相应研究. 我们主要的研究兴趣在于: 对 $f(x,y)$ 加什么条件, 可以保证 $I(x)$ 在区间 $[a,b]$ 上连续、可积以及可导? 下面我们分别来讨论这些问题。

\noindent\textbf{定理 17.1.1}\quad 设函数 $f(x,y)$ 在区域 $D=[a,b]\times[c,d]$ 上连续, 则对于 $\forall x\in[a,b]$, 含参变量定积分 $I(x)=\int_c^d f(x,y)\mathrm dy$ 存在, 并且 $I(x)$ 在区间 $[a,b]$ 上连续。

\noindent\textbf{证明}\quad 由于 $f(x,y)\in C(D)$, 因此对固定的 $x$, $f(x,y)$ 在 $[c,d]$ 上连续, 从而 $I(x)=\int_c^d f(x,y)\mathrm dy$ 存在. 下面证 $I(x)$ 在 $[a,b]$ 上连续. 对于 $\forall x_0\in[a,b]$ 及 $\forall x\in[a,b]$, 有
\[
  I(x)-I(x_0)=\int_c^d [f(x,y)-f(x_0,y)]\mathrm dy.
\]
由 $f(x,y)$ 在 $D$ 上的一致连续性知, 对于 $\forall\varepsilon>0,\exists\delta>0$, 当 $(x_1,y_1),(x_2,y_2)\in D$ 且满足
\[
  \sqrt{(x_1-x_2)^2+(y_1-y_2)^2}<\delta
\]
时, 有
\[
  |f(x_1,y_1)-f(x_2,y_2)|<\frac{\varepsilon}{d-c}.
\]
所以, 当 $|x-x_0|<\delta$ 且 $x\in[a,b]$ 时, 有
\[
  |I(x)-I(x_0)|
  \le \int_c^d |f(x,y)-f(x_0,y)|\mathrm dy
  < \frac{\varepsilon}{d-c}(d-c)=\varepsilon.
\]
证毕。

显然, 定理 17.1.1 中 $f(x,y)$ 在 $D$ 上连续只是 $I(x)$ 连续的充分条件. 读者容易举出 $f(x,y)$ 在 $D$ 上不连续时 $I(x)$ 仍然连续的例子。

另外, 在定理 17.1.1 的条件下, 我们有
\[
  \lim_{x\to x_0}\int_c^d f(x,y)\mathrm dy
  =\lim_{x\to x_0} I(x)
  =I(x_0)
  =\int_c^d \lim_{x\to x_0} f(x,y)\mathrm dy.
\]
因此, 函数 $I(x)$ 的连续性可以解释为极限运算与积分运算的可交换性。

\noindent\textbf{注}\quad 定理 17.1.1 可以作如下推广: 在定理 17.1.1 的假定下, 对于 $\forall(x,u)\in D$, 变上限含参变量积分
\[
  I(x,u)=\int_c^u f(x,y)\mathrm dy
\]
存在, 并且二元函数 $I(x,u)$ 在 $D$ 上连续 (见本章习题). 对于变下限含参变量积分, 也有类似的结论。

由化二重积分为累次积分的定理 (定理 15.3.1), 我们有如下关于含参变量定积分的积分定理。

\noindent\textbf{定理 17.1.2}\quad 设函数 $f(x,y)$ 在区域 $D=[a,b]\times[c,d]$ 上连续, 则函数
\[
  I_1(x)=\int_c^d f(x,y)\mathrm dy
  \quad\text{和}\quad
  I_2(y)=\int_a^b f(x,y)\mathrm dx
\]
分别在区间 $[a,b]$ 和 $[c,d]$ 上可积, 并且
\[
  \int_a^b I_1(x)\mathrm dx=\int_c^d I_2(y)\mathrm dy.
\]
定理 17.1.2 说明, 在定理条件下, 两个求积分的顺序可以交换:
\[
  \int_a^b \mathrm dx\int_c^d f(x,y)\mathrm dy
  =
  \int_c^d \mathrm dy\int_a^b f(x,y)\mathrm dx.
\]

下面我们来讨论含参变量定积分的可导性。

\noindent\textbf{定理 17.1.3}\quad 设函数 $f(x,y)$ 及其偏导数 $f'_x(x,y)$ 在区域 $D=[a,b]\times[c,d]$ 上连续, 则函数 $I(x)=\int_c^d f(x,y)\mathrm dy$ 在区间 $[a,b]$ 上可导, 并且有
\[
  I'(x)=\int_c^d f'_x(x,y)\mathrm dy.
\]

\noindent\textbf{证明}\quad 由于 $f'_x(x,y)$ 在 $D$ 上连续, 从而由定理 17.1.1 知函数
\[
  g(x)=\int_c^d f'_x(x,y)\mathrm dy
\]
在 $[a,b]$ 上存在并且连续. 现将 $g(x)$ 改写成
\[
  g(u)=\int_c^d f'_u(u,y)\mathrm dy.
\]
对于 $\forall x\in[a,b]$, 对 $f'_u(u,y)$ 在 $D=[a,x]\times[c,d]$ 上应用定理 17.1.2, 有
\[
\begin{aligned}
  \int_a^x g(u)\mathrm du
  &=\int_a^x \mathrm du\int_c^d f'_u(u,y)\mathrm dy
    =\int_c^d \mathrm dy\int_a^x f'_u(u,y)\mathrm du\\
  &=\int_c^d [f(x,y)-f(a,y)]\mathrm dy
    = I(x)-I(a).
\end{aligned}
\]
因此对变限定积分 $\int_a^x g(u)\mathrm du$ 求导数得
\[
  I'(x)=g(x)=\int_c^d f'_x(x,y)\mathrm dy.
\]
证毕。

定理 17.1.3 说明了在 $f'_x(x,y)$ 连续的条件下求导数与积分的可交换性. 此定理也可通过导数定义直接求出 $I'(x)$ 来加以证明。

\noindent\textbf{定理 17.1.4}\quad 设函数 $f(x,y)$ 及其偏导数 $f'_x(x,y)$ 在区域 $D=[a,b]\times[c,d]$ 上连续, 且 $\varphi(x)\ (x\in[a,b])$ 是满足 $c\le \varphi(x)\le d$ 的可微函数, 则函数
\[
  I(x)=\int_c^{\varphi(x)} f(x,y)\mathrm dy
\]
在区间 $[a,b]$ 上可导, 并且
\[
  I'(x)=\int_c^{\varphi(x)} f'_x(x,y)\mathrm dy
  + f(x,\varphi(x))\varphi'(x).
\]

\noindent\textbf{证明}\quad 对任意给定的 $x_0\in[a,b]$, 由 $c\le\varphi(x_0)\le d$ 及 $f(x,y)$ 在 $D$ 的连续性知, $f(x_0,y)$ 在 $[c,\varphi(x_0)]$ 上连续, 所以 $I(x_0)=\int_c^{\varphi(x_0)}f(x,y)\mathrm dy$ 存在. 因此 $I(x)$ 在 $[a,b]$ 上有定义。

另外, 令 $u=\varphi(x)$, 则 $I(x)$ 可以看成是 $F(x,u)=\int_c^u f(x,y)\mathrm dy$ 与 $u=\varphi(x)$ 的复合函数. 由定理 17.1.3 我们知
\[
  \frac{\partial F(x,u)}{\partial x}=\int_c^u f'_x(x,y)\mathrm dy.
\]
再由定理 17.1.1 的注知 $\frac{\partial F(x,u)}{\partial x}=\int_c^u f'_x(x,y)\mathrm dy$ 在 $D$ 上连续. 对 $F(x,u)=\int_c^u f(x,y)\mathrm dy$ 两边关于 $u$ 求偏导数得
\[
  \frac{\partial F(x,u)}{\partial u}=f(x,u),
\]
由假设它也在 $D$ 上连续. 这说明 $F(x,u)$ 的两个偏导数在 $D$ 上连续. 因此利用复合函数求导法得
\[
\begin{aligned}
  I'(x)
  &=\frac{\partial F(x,u)}{\partial x}
    +\frac{\partial F(x,u)}{\partial u}\cdot\frac{\partial u}{\partial x}\\
  &=\int_c^{\varphi(x)} f'_x(x,y)\mathrm dy
    +f(x,\varphi(x))\varphi'(x).
\end{aligned}
\]
证毕。

\noindent\textbf{例 17.1.1}\quad 求极限
\[
  \lim_{x\to0+0}\int_0^1 (1+y)^x e^{xy}\sin(x+y)\mathrm dy.
\]

\noindent\textbf{解}\quad 由于 $f(x,y)=(1+y)^x e^{xy}\sin(x+y)$ 在 $[0,+\infty)\times[0,1]$ 上连续, 因此 $I(x)=\int_0^1(1+y)^x e^{xy}\sin(x+y)\mathrm dy$ 在 $[0,+\infty)$ 上连续. 特别地, 我们有
\[
  \lim_{x\to0+0}\int_0^1 (1+y)^x e^{xy}\sin(x+y)\mathrm dy
  =\int_0^1 \sin y\,\mathrm dy
  =1-\cos1.
\]

\noindent\textbf{例 17.1.2}\quad 求定积分
\[
  \int_0^1 \frac{x^{2e-1}-x^{e-1}}{\ln x}\mathrm dx.
\]

\noindent\textbf{解}\quad 如果直接求此定积分, 计算将会非常困难, 因此考虑利用含参变量积分来计算. 由
\[
  \frac{x^{2e-1}-x^{e-1}}{\ln x}
  =
  \int_{e-1}^{2e-1} x^y\mathrm dy
\]
及 $f(x,y)=x^y$ 在 $[0,1]\times[e-1,2e-1]$ 上连续, 我们有
\[
\begin{aligned}
  \int_0^1 \frac{x^{2e-1}-x^{e-1}}{\ln x}\mathrm dx
  &=\int_0^1 \mathrm dx\int_{e-1}^{2e-1}x^y\mathrm dy\\
  &=\int_{e-1}^{2e-1}\mathrm dy\int_0^1x^y\mathrm dx
    =\int_{e-1}^{2e-1}\frac{\mathrm dy}{1+y}\\
  &=\left.\ln(1+y)\right|_{e-1}^{2e-1}=\ln2.
\end{aligned}
\]

\noindent\textbf{例 17.1.3}\quad 设函数 $f(s,t)$ 在 $\mathbb R^2$ 上连续, 记
\[
  F(x)=\int_x^{x^2}\mathrm ds\int_s^x f(s,t)\mathrm dt,
\]

\noindent 求 $F'(x)$。

\noindent\textbf{解}\quad 记 $g(s,x)=\int_s^x f(s,t)\mathrm dt$. 由定理 17.1.1 的注知, 它在 $\mathbb R^2$ 中连续, 从而对于 $\forall x\in(-\infty,+\infty)$, 我们有
\[
  F(x)=\int_x^{x^2}g(s,x)\mathrm ds
  =\int_0^{x^2}g(s,x)\mathrm ds-\int_0^x g(s,x)\mathrm ds
\]
存在. 另外, $\frac{\partial g(s,x)}{\partial x}=f(s,x)$ 在 $\mathbb R^2$ 上连续, 并且 $x,x^2$ 均是可导函数, 从而由定理 17.1.4 知 $F'(x)$ 存在。

记
\[
  F_1(x)=\int_0^{x^2}g(s,x)\mathrm ds,\qquad
  F_2(x)=\int_0^x g(s,x)\mathrm ds,
\]
则有
\[
\begin{aligned}
  F_1'(x)
  &=\frac{\mathrm d}{\mathrm dx}\left(\int_0^{x^2}g(s,x)\mathrm ds\right)
    =\int_0^{x^2}\frac{\partial g(s,x)}{\partial x}\mathrm ds
      +g(x^2,x)(x^2)'\\
  &=\int_0^{x^2}f(s,x)\mathrm ds
    +2x\int_{x^2}^x f(x^2,t)\mathrm dt.
\end{aligned}
\]
同理,
\[
  F_2'(x)=\int_0^x f(s,x)\mathrm ds+\int_x^x f(x,t)\mathrm dt
  =\int_0^x f(s,x)\mathrm ds.
\]
最后我们有
\[
  F'(x)=F_1'(x)-F_2'(x)
  =\int_x^{x^2}f(s,x)\mathrm ds+2x\int_{x^2}^x f(x^2,t)\mathrm dt.
\]

\section{含参变量广义积分}

含参变量广义积分视积分为无穷积分或瑕积分可分为两种: 含参变量无穷积分和含参变量瑕积分. 我们先来讨论含参变量无穷积分。

\subsection{含参变量无穷积分}

设函数 $f(x,y)$ 在 $E\times[c,+\infty)$ 上有定义, 其中 $E\subset\mathbb R$ 为一个集合. 若对于 $\forall x\in E$, 广义积分 $\int_c^{+\infty}f(x,y)\mathrm dy$ 收敛, 则可得到 $E$ 上的函数
\[
  I(x)=\int_c^{+\infty}f(x,y)\mathrm dy,
\]
我们称该函数为\textbf{含参变量无穷积分}。

在这里我们仍将讨论 $I(x)$ 的连续性、可微性以及积分与无穷积分的交换性等问题. 在讨论这些问题时, 一般我们假定 $E$ 是一个区间. 如同函数项级数, 为了保证 $I(x)$ 能够继承 $f(x,y)$ 的一些性质, 引入含参变量无穷积分一致收敛的概念是必要的。

\noindent\textbf{定义 17.2.1}\quad 设函数 $f(x,y)$ 在 $E\times[c,+\infty)$ 上有定义, 其中 $E\subset\mathbb R$ 是一个区间. 若对于 $\forall\varepsilon>0,\exists A_0>0$, 当 $A>A_0$ 时, 对于 $\forall x\in E$, 有
\[
  \left|\int_A^{+\infty}f(x,y)\mathrm dy\right|<\varepsilon,
\]
则称含参变量无穷积分 $\int_c^{+\infty}f(x,y)\mathrm dy$ 在 $E$ 上\textbf{一致收敛}。

在一致收敛的研究中, 柯西准则是一个非常重要的理论工具。

\noindent\textbf{定理 17.2.1}\quad 设函数 $f(x,y)$ 在 $E\times[c,+\infty)$ 上有定义, 其中 $E\subset\mathbb R$ 是一个区间, 则含参变量无穷积分 $\int_c^{+\infty}f(x,y)\mathrm dy$ 在 $E$ 上一致收敛的充分必要条件是: 对于 $\forall\varepsilon>0,\exists A_0>0$, 当 $A,A'>A_0$ 时, 对于 $\forall x\in E$ 有
\[
  \left|\int_A^{A'}f(x,y)\mathrm dy\right|<\varepsilon.
\]

\noindent\textbf{证明}\quad 必要性\quad 设 $\int_c^{+\infty}f(x,y)\mathrm dy$ 在 $E$ 上一致收敛, 则对于 $\forall\varepsilon>0,\exists A_0>0$, 当 $A>A_0$ 时, 对于 $\forall x\in E$ 有
\[
  \left|\int_A^{+\infty}f(x,y)\mathrm dy\right|<\frac{\varepsilon}{2};
\]
从而当 $A,A'>A_0$ 时, 有
\[
  \left|\int_A^{A'}f(x,y)\mathrm dy\right|
  \le\left|\int_A^{+\infty}f(x,y)\mathrm dy\right|
    +\left|\int_{A'}^{+\infty}f(x,y)\mathrm dy\right|
  <\frac{\varepsilon}{2}+\frac{\varepsilon}{2}=\varepsilon.
\]

充分性\quad 设对于 $\forall\varepsilon>0,\exists A_0>0$, 当 $A,A'>A_0$ 时, 对于 $\forall x\in E$ 有
\[
  \left|\int_A^{A'}f(x,y)\mathrm dy\right|<\frac{\varepsilon}{2}.
  \tag{17.2.1}
\]
这说明 $\int_c^{+\infty}f(x,y)\mathrm dy$ 在每点 $x\in E$ 满足柯西准则, 从而它点点收敛. 在式 (17.2.1) 中令 $A'\to+\infty$, 则对于 $\forall x\in E$, 都有
\[
  \left|\int_A^{+\infty}f(x,y)\mathrm dy\right|\le\frac{\varepsilon}{2}<\varepsilon.
\]
证毕。

设函数 $f(x,y)$ 在 $E\times[c,+\infty)$ 上有定义, 其中 $E\subset\mathbb R$ 是一个区间. 若对于 $\forall x\in E$, $\int_c^{+\infty}|f(x,y)|\mathrm dy$ 收敛, 则称 $\int_c^{+\infty}f(x,y)\mathrm dy$ 在 $E$ 上\textbf{绝对收敛}. 显然, 若 $\int_c^{+\infty}f(x,y)\mathrm dy$ 在 $E$ 上绝对收敛, 则 $\int_c^{+\infty}f(x,y)\mathrm dy$ 在 $E$ 上收敛. 另外, 若 $\int_c^{+\infty}|f(x,y)|\mathrm dy$ 在 $E$ 上一致收敛, 则称 $\int_c^{+\infty}f(x,y)\mathrm dy$ 在 $E$ 上\textbf{绝对一致收敛}。

\noindent\textbf{定理 17.2.2（魏尔斯特拉斯定理）}\quad 设函数 $f(x,y)$ 在 $E\times[c,+\infty)$ 上有定义, 其中 $E\subset\mathbb R$ 是一个区间. 若存在函数 $g(y)(y\in[c,+\infty))$, 使得对于 $\forall x\in E$ 及 $\forall y\in[c,+\infty)$, 有 $|f(x,y)|\le g(y)$, 并且 $\int_c^{+\infty}g(y)\mathrm dy$ 收敛, 则 $\int_c^{+\infty}f(x,y)\mathrm dy$ 在 $E$ 上绝对一致收敛。

\noindent\textbf{证明}\quad 对于 $\forall\varepsilon>0$, 由 $\int_c^{+\infty}g(y)\mathrm dy$ 收敛, 从而 $\exists A_0>0$, 当 $A,A'>A_0$ 时, 有
\[
  \left|\int_A^{A'}g(y)\mathrm dy\right|<\varepsilon.
\]
我们不妨设 $A'>A$, 从而对于 $\forall x\in E$, 有
\[
  \left|\int_A^{A'}f(x,y)\mathrm dy\right|
  \le\int_A^{A'}|f(x,y)|\mathrm dy
  \le\int_A^{A'}g(y)\mathrm dy<\varepsilon.
\]
由定理 17.2.1 知 $\int_c^{+\infty}|f(x,y)|\mathrm dy$ 在 $E$ 上一致收敛, 即 $\int_c^{+\infty}f(x,y)\mathrm dy$ 在 $E$ 上绝对一致收敛. 证毕。

对于含参变量无穷积分一致收敛的判别法则, 可以参照无穷积分的收敛判别法则来加以修改得到. 事实上, 只要将后者的判别法则中的条件对参数一致满足即可得到前者的判别法则. 我们有下面的两个定理。

\noindent\textbf{定理 17.2.3（狄利克雷判别法）}\quad 设函数 $f(x,y),g(x,y)$ 在 $E\times[c,+\infty)$ 上有定义 (其中 $E\subset\mathbb R$ 是一个区间), 并且满足:
\[
\begin{gathered}
  (1)\quad \text{存在 } M>0,\ \text{对于 } \forall A>c \text{ 及 } \forall x\in E,\ \text{有 }
  \left|\int_c^A f(x,y)\mathrm dy\right|\le M;\\
  (2)\quad \text{对任意固定的 } x\in E,\ g(x,y) \text{ 是 } y \text{ 的单调函数, 且对于 } \forall\varepsilon>0,\exists A>c,\\
  \text{当 } y>A \text{ 时, 对一切 } x\in E,\ \text{有 } |g(x,y)|<\varepsilon,
  \text{ 即当 } y\to+\infty \text{ 时, } g(x,y) \text{ 关于 } x \text{ 一致趋于 } 0,
\end{gathered}
\]
则含参变量无穷积分 $\int_c^{+\infty}f(x,y)g(x,y)\mathrm dy$ 在 $E$ 上一致收敛。

\noindent\textbf{证明}\quad 对于 $A'>A>c$, 由定积分第二中值定理, $\exists\xi\in(A,A')$, 使得
\[
  \int_A^{A'}f(x,y)g(x,y)\mathrm dy
  =
  g(x,A)\int_A^\xi f(x,y)\mathrm dy
  +g(x,A')\int_\xi^{A'}f(x,y)\mathrm dy.
  \tag{17.2.2}
\]
由 (2) 知, 对于 $\forall\varepsilon>0,\exists A_0>c$, 当 $y>A_0$ 时, 对一切 $x\in E$, 有 $|g(x,y)|<\frac{\varepsilon}{4M}$. 因此当 $A'>A>A_0$ 时, 由式 (17.2.2) 即有
\[
\begin{aligned}
  \left|\int_A^{A'}f(x,y)g(x,y)\mathrm dy\right|
  &<\frac{\varepsilon}{4M}\left|\int_c^\xi f(x,y)\mathrm dy-\int_c^A f(x,y)\mathrm dy\right|\\
  &\quad+\frac{\varepsilon}{4M}\left|\int_c^{A'}f(x,y)\mathrm dy-\int_c^\xi f(x,y)\mathrm dy\right|\\
  &\le\frac{\varepsilon}{4M}(2M+2M)=\varepsilon.
\end{aligned}
\]
证毕。

\noindent\textbf{定理 17.2.4（阿贝尔判别法）}\quad 设函数 $f(x,y),g(x,y)$ 在 $E\times[c,+\infty)$ 上有定义 (其中 $E\subset\mathbb R$ 是一个区间), 并且满足:
\[
\begin{gathered}
  (1)\quad \int_c^{+\infty}f(x,y)\mathrm dy \text{ 在 } E \text{ 上一致收敛};\\
  (2)\quad \text{对任意固定的 } x\in E,\ g(x,y) \text{ 是 } y \text{ 的单调函数, 并且存在常数 } M>0,\\
  \text{对于 } \forall x\in E \text{ 及 } \forall y\in[c,+\infty),\ \text{有 } |g(x,y)|\le M,
\end{gathered}
\]
则含参变量无穷积分 $\int_c^{+\infty}f(x,y)g(x,y)\mathrm dy$ 在 $E$ 上一致收敛。

\noindent\textbf{证明}\quad 对于 $A'>A>c$, 我们仍然有式 (17.2.2)。

由 $\int_c^{+\infty}f(x,y)\mathrm dy$ 在 $E$ 上一致收敛, 对于 $\forall\varepsilon>0,\exists A_0>c$, 当 $A'>A>A_0$ 时, 对一切 $x\in E$, 有
\[
  \left|\int_A^{A'}f(x,y)\mathrm dy\right|<\frac{\varepsilon}{2M},
\]
从而由式 (17.2.2) 得
\[
  \left|\int_A^{A'}f(x,y)g(x,y)\mathrm dy\right|
  <\frac{\varepsilon}{2M}\bigl(|g(x,A)|+|g(x,A')|\bigr)\le\varepsilon.
\]
证毕。

\noindent\textbf{例 17.2.1}\quad 证明含参变量无穷积分 $\int_0^{+\infty}y^xe^{-y}\mathrm dy$, 满足:
\[
\begin{gathered}
  (1)\quad \text{在 } [0,M]\ (0<M<+\infty) \text{ 上一致收敛};\\
  (2)\quad \text{在 } [0,+\infty) \text{ 上不一致收敛}.
\end{gathered}
\]

\noindent\textbf{证明}\quad (1) 对于 $\forall x\in[0,M]$, 有
\[
  0\le y^xe^{-y}\le y^M e^{-y}.
\]
由于 $\int_0^{+\infty}y^M e^{-y}\mathrm dy$ 收敛, 由定理 17.2.4 知, $\int_0^{+\infty}y^xe^{-y}\mathrm dy$ 在 $[0,M]$ 上一致收敛。

(2) 对任意的 $A_0>1$, 任取 $A_1>A_0,A_2=A_1+1$, 考查 $\int_{A_1}^{A_2}y^xe^{-y}\mathrm dy$. 由于
\[
  \lim_{x\to+\infty}A_1^x e^{-A_2}=+\infty,
\]
从而存在 $x'\in(0,+\infty)$, 使得 $A_1^{x'}e^{-A_2}\ge1$, 因此
\[
  \int_{A_1}^{A_2}y^{x'}e^{-y}\mathrm dy
  \ge \int_{A_1}^{A_2}A_1^{x'}e^{-A_2}\mathrm dy\ge1.
\]
这说明 $\int_0^{+\infty}y^xe^{-y}\mathrm dy$ 在 $[0,+\infty)$ 上不一致收敛。

\noindent\textbf{例 17.2.2}\quad 证明含参变量无穷积分 $\int_0^{+\infty}\frac{\sin xy}{y}\mathrm dy$
\[
\begin{gathered}
  (1)\quad \text{在 } [\alpha_0,+\infty)\ (\alpha_0>0) \text{ 上一致收敛};\\
  (2)\quad \text{在 } (0,+\infty) \text{ 内不一致收敛}.
\end{gathered}
\]

\noindent\textbf{证明}\quad (1) 方法 1\quad 设 $f(x,y)=\sin xy,g(x,y)=\frac1y$, 则对于 $\forall x\in[\alpha_0,+\infty)$ 及 $\forall A>0$, 有
\[
  \left|\int_0^A f(x,y)\mathrm dy\right|
  =
  \left|\int_0^A\sin xy\,\mathrm dy\right|
  =
  \left|\frac1x(1-\cos xA)\right|
  \le\frac2{\alpha_0}.
\]
而 $g(x,y)=\frac1y$ 关于 $y$ 单调, 且对 $x\in[\alpha_0,+\infty)$ 一致趋于 $0$, 由狄利克雷判别法知 $\int_0^{+\infty}\frac{\sin xy}{y}\mathrm dy$ 在 $[\alpha_0,+\infty)$ 上一致收敛。

方法 2\quad 对于 $\forall\varepsilon>0$, 由 $\int_0^{+\infty}\frac{\sin y}{y}\mathrm dy$ 收敛, 从而存在 $A_0>0$, 使得当 $A>A_0$ 时, 有
\[
  \left|\int_A^{+\infty}\frac{\sin y}{y}\mathrm dy\right|<\varepsilon.
\]
取 $A_1=\frac{A_0}{\alpha_0}$, 则当 $A>A_1$ 时, 对一切 $x\in[\alpha_0,+\infty)$, 有 $xA>\alpha_0 A_1=A_0$, 从而有
\[
  \left|\int_A^{+\infty}\frac{\sin xy}{y}\mathrm dy\right|
  =
  \left|\int_{xA}^{+\infty}\frac{\sin y}{y}\mathrm dy\right|
  <\varepsilon.
\]
因此 $\int_0^{+\infty}\frac{\sin xy}{y}\mathrm dy$ 在 $[\alpha_0,+\infty)$ 上一致收敛。

(2) 取 $x=\frac1{2k}\ (k\in\mathbb N)$, 我们有
\[
\begin{aligned}
  \left|\int_{4k\pi}^{5k\pi}\frac{\sin xy}{y}\mathrm dy\right|
  &=\left|\int_{4k\pi}^{5k\pi}\frac{\sin\frac{y}{2k}}{y}\mathrm dy\right|\\
  &\ge \frac1{5k\pi}\left|\int_{4k\pi}^{5k\pi}\sin\frac{y}{2k}\mathrm dy\right|
  =\frac2{5\pi}.
\end{aligned}
\]
因此存在 $\varepsilon_0=\frac1{5\pi}$, 对 $\forall A>0$, 存在正整数 $k_0>A$, 从而存在 $A'=4k_0\pi,A''=5k_0\pi>A$ 及 $x'=\frac1{2k_0}\in(0,+\infty)$, 使得
\[
  \left|\int_{A'}^{A''}\frac{\sin x'y}{y}\mathrm dy\right|>\varepsilon_0.
\]
这说明 $\int_0^{+\infty}\frac{\sin xy}{y}\mathrm dy$ 在 $(0,+\infty)$ 内不一致收敛。

\noindent\textbf{例 17.2.3}\quad 设函数 $f(x)$ 在 $[0,+\infty)$ 内连续, 证明无穷积分
\[
  \int_0^{+\infty}e^{-\alpha x}f(x)\mathrm dx
\]
对 $\alpha\in(0,+\infty)$ 一致收敛的充分必要条件是无穷积分 $\int_0^{+\infty}f(x)\mathrm dx$ 收敛。

\noindent\textbf{证明}\quad 充分性. 若无穷积分 $\int_0^{+\infty}f(x)\mathrm dx$ 收敛, 则它对参数 $\alpha\in(0,+\infty)$ 一致收敛. 又 $e^{-\alpha x}$ 关于 $x$ 单调, 且对 $\alpha\in(0,+\infty)$ 有 $|e^{-\alpha x}|\le1$, 由阿贝尔判别法知 $\int_0^{+\infty}e^{-\alpha x}f(x)\mathrm dx$ 在 $(0,+\infty)$ 上一致收敛。

必要性. 假设 $\int_0^{+\infty}f(x)\mathrm dx$ 不收敛, 则存在 $\varepsilon_0>0$, 对任意 $A_0>0$, 存在 $A_2>A_1>A_0$, 使得
\[
  \left|\int_{A_1}^{A_2}f(x)\mathrm dx\right|>2\varepsilon_0.
\]
由于 $e^{-\alpha x}f(x)$ 在 $[A_1,A_2]\times[0,+\infty)$ 上连续, 所以
\[
  \lim_{\alpha\to0+0}\int_{A_1}^{A_2}e^{-\alpha x}f(x)\mathrm dx
  =\int_{A_1}^{A_2}f(x)\mathrm dx.
\]
于是存在 $\alpha'>0$, 使得
\[
  \left|\int_{A_1}^{A_2}e^{-\alpha'x}f(x)\mathrm dx\right|
  \ge \frac12\left|\int_{A_1}^{A_2}f(x)\mathrm dx\right|
  >\varepsilon_0.
\]
这说明 $\int_0^{+\infty}e^{-\alpha x}f(x)\mathrm dx$ 在 $(0,+\infty)$ 上不一致收敛, 矛盾. 证毕。

\subsection{含参变量无穷积分的性质}

本节讨论含参变量无穷积分的连续性、可积性与可微性。

\noindent\textbf{定理 17.2.5}\quad 设函数 $f(x,y)$ 在 $E\times[c,+\infty)$ 上有定义, 其中 $E\subset\mathbb R$, 则含参变量无穷积分
\[
  \int_c^{+\infty}f(x,y)\mathrm dy
\]
在 $E$ 上一致收敛的充分必要条件是: 对任意的满足条件
\[
  c<t_1<t_2<\cdots<t_k<\cdots,\qquad \lim_{k\to\infty}t_k=+\infty
\]
的序列 $\{t_k\}$, 函数列
\[
  F_k(x)=\int_c^{t_k}f(x,y)\mathrm dy\quad(k=1,2,\cdots)
\]
在 $E$ 上一致收敛。

\noindent\textbf{证明}\quad 必要性. 设 $\int_c^{+\infty}f(x,y)\mathrm dy$ 在 $E$ 上一致收敛. 对 $\forall\varepsilon>0$, 存在 $A_0>0$, 使得当 $A',A''>A_0$ 时, 对一切 $x\in E$, 有
\[
  \left|\int_{A'}^{A''}f(x,y)\mathrm dy\right|<\varepsilon.
\]
任取满足条件的序列 $\{t_k\}$, 当 $k,k'$ 充分大时, 有 $t_k,t_{k'}>A_0$, 故
\[
  |F_k(x)-F_{k'}(x)|
  =\left|\int_{t_{k'}}^{t_k}f(x,y)\mathrm dy\right|<\varepsilon
  \quad(x\in E),
\]
即 $\{F_k(x)\}$ 在 $E$ 上一致收敛。

充分性. 反设 $\int_c^{+\infty}f(x,y)\mathrm dy$ 在 $E$ 上不一致收敛. 则存在 $\varepsilon_0>0$, 对任意 $A>0$, 存在 $A',A''>A$ 及 $x'\in E$, 使得
\[
  \left|\int_{A'}^{A''}f(x',y)\mathrm dy\right|\ge\varepsilon_0.
\]
取 $A=1$, 得到 $t_1',t_1''$ 与 $x_1\in E$; 依次取 $A=t_k''$, 可得到序列
\[
  k<t_k'<t_k''<t_{k+1}'<t_{k+1}''\quad(k=1,2,\cdots)
\]
及 $x_k\in E$, 满足
\[
  \left|\int_{t_k'}^{t_k''}f(x_k,y)\mathrm dy\right|\ge\varepsilon_0.
\]
令 $\{t_k\}=\{t_1',t_1'',t_2',t_2'',\cdots\}$, 则相应的函数列 $F_k(x)=\int_c^{t_k}f(x,y)\mathrm dy$ 不在 $E$ 上一致收敛, 矛盾. 证毕。

容易看出, 在定理 17.2.5 中, 若 $\int_c^{+\infty}f(x,y)\mathrm dy$ 在 $E$ 上一致收敛, 则对任意满足该定理条件的 $\{t_k\}$, 当 $k\to\infty$ 时, 都有
\[
  F_k(x)\rightrightarrows \int_c^{+\infty}f(x,y)\mathrm dy\quad(x\in E).
\]

\noindent\textbf{定理 17.2.6}\quad 设函数 $f(x,y)$ 在 $E\times[c,+\infty)$ 上连续, 其中 $E\subset\mathbb R$ 是一个区间, 并且含参变量无穷积分 $\int_c^{+\infty}f(x,y)\mathrm dy$ 在 $E$ 上一致收敛到函数 $I(x)$, 则 $I(x)$ 在 $E$ 上连续。

\noindent\textbf{证明}\quad 取 $t_k=c+k\ (k=1,2,\cdots)$, 令
\[
  F_k(x)=\int_c^{t_k}f(x,y)\mathrm dy.
\]
由定理 17.1.1 知, $F_k(x)$ 在 $E$ 上连续. 又由定理 17.2.5 知 $F_k(x)\rightrightarrows I(x)$, 因此 $I(x)$ 在 $E$ 上连续。

\noindent\textbf{定理 17.2.7}\quad 设函数 $f(x,y)$ 在 $[a,b]\times[c,+\infty)$ 上连续, 且含参变量无穷积分 $\int_c^{+\infty}f(x,y)\mathrm dy$ 在 $[a,b]$ 上一致收敛, 则有
\begin{equation}
  \int_a^b\mathrm dx\int_c^{+\infty}f(x,y)\mathrm dy
  =
  \int_c^{+\infty}\mathrm dy\int_a^b f(x,y)\mathrm dx.
  \tag{17.2.3}
\end{equation}

\noindent\textbf{证明}\quad 由定理 17.2.6 知 $I(x)=\int_c^{+\infty}f(x,y)\mathrm dy$ 在 $[a,b]$ 上连续. 对任意满足定理 17.2.5 条件的序列 $\{t_k\}$, 令 $F_k(x)=\int_c^{t_k}f(x,y)\mathrm dy$, 则 $F_k(x)\rightrightarrows I(x)$. 因此
\[
\begin{aligned}
  \int_a^b I(x)\mathrm dx
  &=\int_a^b\lim_{k\to\infty}F_k(x)\mathrm dx\\
  &=\lim_{k\to\infty}\int_a^b\mathrm dx\int_c^{t_k}f(x,y)\mathrm dy\\
  &=\lim_{k\to\infty}\int_c^{t_k}\mathrm dy\int_a^b f(x,y)\mathrm dx.
\end{aligned}
\]
由于 $\{t_k\}$ 是任意的, 故 $\int_c^{+\infty}\mathrm dy\int_a^b f(x,y)\mathrm dx$ 收敛且等式 (17.2.3) 成立。

\noindent\textbf{定理 17.2.8}\quad 设函数 $f(x,y)$ 及其偏导数 $f'_x(x,y)$ 在 $E\times[c,+\infty)$ 上连续, 其中 $E\subset\mathbb R$ 是一个区间, 再设存在 $x_0\in E$, 使得 $\int_c^{+\infty}f(x_0,y)\mathrm dy$ 收敛, 并且 $\int_c^{+\infty}f'_x(x,y)\mathrm dy$ 在 $E$ 上一致收敛, 则
\[
\begin{gathered}
  (1)\quad I(x)=\int_c^{+\infty}f(x,y)\mathrm dy \text{ 在 } E \text{ 上一致收敛};\\
  (2)\quad I'(x)=\left(\int_c^{+\infty}f(x,y)\mathrm dy\right)'
  =\int_c^{+\infty}f'_x(x,y)\mathrm dy.
\end{gathered}
\]

\noindent\textbf{证明}\quad 任取满足定理 17.2.5 条件的序列 $\{t_k\}$, 令
\[
  F_k(x)=\int_c^{t_k}f(x,y)\mathrm dy.
\]
由定理 17.1.3 知
\[
  F'_k(x)=\int_c^{t_k}f'_x(x,y)\mathrm dy.
\]
因为数列 $\{F_k(x_0)\}$ 收敛, 且函数列 $\{F'_k(x)\}$ 在 $E$ 上一致收敛, 所以 $\{F_k(x)\}$ 在 $E$ 上一致收敛, 并且可逐项求导. 于是结论成立。

\noindent\textbf{例 17.2.5}\quad 计算狄利克雷积分 $\int_0^{+\infty}\frac{\sin x}{x}\mathrm dx$。

\noindent\textbf{解}\quad 令
\[
  I(\alpha)=\int_0^{+\infty}e^{-\alpha x}\frac{\sin x}{x}\mathrm dx\quad(\alpha>0).
\]
由阿贝尔判别法知该积分对 $\alpha\in(0,+\infty)$ 一致收敛. 对任意 $\alpha_0>0$, $\int_0^{+\infty}(-e^{-\alpha x}\sin x)\mathrm dx$ 在 $[\alpha_0,+\infty)$ 上一致收敛, 故
\[
\begin{aligned}
  I'(\alpha)
  &=-\int_0^{+\infty}e^{-\alpha x}\sin x\,\mathrm dx\\
  &=e^{-\alpha x}\cos x\Big|_0^{+\infty}
    +\alpha\int_0^{+\infty}e^{-\alpha x}\cos x\,\mathrm dx\\
  &=-1+\alpha e^{-\alpha x}\sin x\Big|_0^{+\infty}
    +\alpha^2\int_0^{+\infty}e^{-\alpha x}\sin x\,\mathrm dx.
\end{aligned}
\]
从而 $I'(\alpha)=-\frac1{1+\alpha^2}$, 所以
\[
  I(\alpha)=-\arctan\alpha+C.
\]
又
\[
  |I(\alpha)|\le\int_0^{+\infty}e^{-\alpha x}\mathrm dx=\frac1\alpha,
\]
故 $\lim_{\alpha\to+\infty}I(\alpha)=0$, 从而 $C=\frac\pi2$. 由于 $I(\alpha)$ 在 $\alpha=0$ 处右连续, 得
\[
  \int_0^{+\infty}\frac{\sin x}{x}\mathrm dx
  =\lim_{\alpha\to0+0}I(\alpha)=\frac\pi2.
\]

\subsection{含参变量瑕积分}

设函数 $f(x,y)$ 在 $[a,b]\times(c,d]$ 上连续, 当 $x\in[a,b]$ 时, $f(x,y)$ 以 $c$ 为瑕点. 若对任意 $x\in[a,b]$, 瑕积分
\begin{equation}
  I(x)=\int_c^d f(x,y)\mathrm dy
  \tag{17.2.4}
\end{equation}
收敛, 则 $I(x)$ 为在 $[a,b]$ 上有定义的函数. 通常称 $I(x)$ 为含参变量瑕积分, $x$ 为参变量。

类似于含参变量无穷积分, 可定义含参变量瑕积分的一致收敛性, 并建立柯西准则、魏尔斯特拉斯判别法、狄利克雷判别法、阿贝尔判别法以及连续性、可积性和可微性等结论. 例如作变量替换
\[
  y=c+\frac1t,
\]
则
\[
  \int_c^d f(x,y)\mathrm dy
  =
  \int_{\frac1{d-c}}^{+\infty}
  f\left(x,c+\frac1t\right)\frac{\mathrm dt}{t^2},
\]
从而含参变量瑕积分可化为含参变量无穷积分来讨论。

上述结论也可以推广到含多个参变量的积分. 例如, 设 $S$ 为空间中的一片光滑曲面, $\mathbf F(x,y,z,t)$ 为流体在点 $(x,y,z)$、时刻 $t$ 的速度场, $\mathbf n(x,y,z)$ 为曲面的单位法向量. 则单位时间内通过曲面 $S$ 的流量为
\[
  S(t)=\iint_S \mathbf F(x,y,z,t)\cdot\mathbf n(x,y,z)\mathrm dS.
\]
若 $\mathbf F$ 关于 $t$ 连续, 则 $S(t)$ 在相应区间上连续. 因而在时间区间 $[t_0,t_1]$ 内通过曲面 $S$ 的流量为
\[
  Q=\int_{t_0}^{t_1}\mathrm dt
  \iint_S \mathbf F(x,y,z,t)\cdot\mathbf n(x,y,z)\mathrm dS.
\]
特别地, 若速度场与时间无关, 则
\[
  Q=\int_{t_0}^{t_0+1}\mathrm dt
  \iint_S \mathbf F(x,y,z)\cdot\mathbf n(x,y,z)\mathrm dS
  =\iint_S \mathbf F(x,y,z)\cdot\mathbf n(x,y,z)\mathrm dS.
\]

\section{$\Gamma$ 函数与 B 函数}

\subsection{$\Gamma$ 函数}

\noindent\textbf{定义 17.3.1}\quad 称含参变量积分
\[
  \Gamma(s)=\int_0^{+\infty}x^{s-1}e^{-x}\mathrm dx
\]
为 $\Gamma$ 函数。

当 $s<1$ 时, $x=0$ 为瑕点. 将积分分成
\[
  \int_0^{+\infty}x^{s-1}e^{-x}\mathrm dx
  =
  \int_0^1x^{s-1}e^{-x}\mathrm dx+\int_1^{+\infty}x^{s-1}e^{-x}\mathrm dx.
\]
第一个积分当且仅当 $s>0$ 时收敛, 第二个积分对任意实数 $s$ 都收敛. 因此 $\Gamma$ 函数的定义域为 $(0,+\infty)$, 且 $\Gamma(s)>0$。

\noindent\textbf{性质 17.3.1}\quad $\Gamma(s+1)=s\Gamma(s)\ (s>0)$. 特别地, 当 $k$ 为非负整数时,
\[
  \Gamma(k+1)=k!.
\]

\noindent\textbf{证明}\quad
\[
\begin{aligned}
  \Gamma(s+1)
  &=\int_0^{+\infty}x^s e^{-x}\mathrm dx
  =-\int_0^{+\infty}x^s\mathrm d e^{-x}\\
  &=-x^s e^{-x}\Big|_0^{+\infty}
  +s\int_0^{+\infty}x^{s-1}e^{-x}\mathrm dx
  =s\Gamma(s).
\end{aligned}
\]
由此递推即得 $\Gamma(k+1)=k!$。

\noindent\textbf{性质 17.3.2}\quad 对 $s>0$, 有
\[
  \Gamma(s)=2\int_0^{+\infty}x^{2s-1}e^{-x^2}\mathrm dx.
\]

\noindent\textbf{证明}\quad 在 $\Gamma(s)$ 的定义式中令 $x=t^2$, 即得。

\noindent\textbf{性质 17.3.3}\quad $\Gamma(s)\in C^\infty(0,+\infty)$。

\noindent\textbf{证明}\quad 对任意非负整数 $k$,
\[
  \frac{\partial^k}{\partial s^k}\bigl(x^{s-1}e^{-x}\bigr)
  =x^{s-1}(\ln x)^k e^{-x}.
\]
任取 $s_0>0$. 当 $0<x\le1$ 且 $s\in[\frac{s_0}2,s_0+1]$ 时,
\[
  x^{s-1}|\ln x|^k e^{-x}\le x^{\frac{s_0}2-1}|\ln x|^k,
\]
右端在 $(0,1]$ 上可积; 当 $x\ge1$ 时,
\[
  x^{s-1}|\ln x|^k e^{-x}\le x^{s_0+k}e^{-x},
\]
右端在 $[1,+\infty)$ 上可积. 因此
\[
  \int_0^{+\infty}x^{s-1}(\ln x)^k e^{-x}\mathrm dx
\]
在 $[\frac{s_0}2,s_0+1]$ 上一致收敛, 从而
\[
  \Gamma^{(k)}(s_0)=\int_0^{+\infty}x^{s_0-1}(\ln x)^k e^{-x}\mathrm dx.
\]
由于 $s_0$ 与 $k$ 任意, 得 $\Gamma(s)\in C^\infty(0,+\infty)$。

\noindent\textbf{性质 17.3.4}\quad $\Gamma(s)$ 与 $\ln\Gamma(s)$ 在 $(0,+\infty)$ 内都是严格凸函数。

\noindent\textbf{证明}\quad 由性质 17.3.3,
\[
  \Gamma''(s)=\int_0^{+\infty}x^{s-1}(\ln x)^2e^{-x}\mathrm dx>0,
\]
故 $\Gamma(s)$ 严格凸. 又
\[
  (\ln\Gamma(s))''=\frac{\Gamma(s)\Gamma''(s)-(\Gamma'(s))^2}{\Gamma^2(s)}.
\]
由柯西不等式,
\[
\begin{aligned}
  \Gamma(s)\Gamma''(s)
  &=
  \int_0^{+\infty}x^{s-1}e^{-x}\mathrm dx
  \int_0^{+\infty}x^{s-1}(\ln x)^2e^{-x}\mathrm dx\\
  &>
  \left(\int_0^{+\infty}|x^{s-1}\ln x\,e^{-x}|\mathrm dx\right)^2
  \ge(\Gamma'(s))^2,
\end{aligned}
\]
因此 $\ln\Gamma(s)$ 严格凸。

\subsection{B 函数}

\noindent\textbf{定义 17.3.2}\quad 称含参变量积分
\[
  B(p,q)=\int_0^1x^{p-1}(1-x)^{q-1}\mathrm dx
\]
为 B 函数, 其中 $p>0,q>0$。

\noindent\textbf{性质 17.3.5}\quad B 函数具有以下性质:
\[
\begin{gathered}
  (1)\quad B(p,q)=B(q,p);\\
  (2)\quad B(p,q)=\frac{p-1}{p+q-1}B(p-1,q)\quad(p>1,q>0).
\end{gathered}
\]

\noindent\textbf{证明}\quad (1) 在定义式中令 $x=1-t$ 即得。

(2) 由分部积分,
\[
\begin{aligned}
  B(p,q)
  &=\int_0^1x^{p-1}(1-x)^{q-1}\mathrm dx\\
  &=-\frac1q\int_0^1x^{p-1}\mathrm d(1-x)^q\\
  &=\frac{p-1}{q}\int_0^1x^{p-2}(1-x)^q\mathrm dx\\
  &=\frac{p-1}{q}\bigl(B(p-1,q)-B(p,q)\bigr).
\end{aligned}
\]
整理即得结论. 特别地, 当 $p>1,q>1$ 时,
\[
  B(p,q)=\frac{(p-1)(q-1)}{(p+q-1)(p+q-2)}B(p-1,q-1).
\]

\noindent\textbf{性质 17.3.6}\quad B 函数具有以下形式:
\[
\begin{gathered}
  (1)\quad B(p,q)=2\int_0^{\frac\pi2}\cos^{2p-1}\theta\sin^{2q-1}\theta\,\mathrm d\theta\quad(p,q>0);\\
  (2)\quad B(p,q)=\int_0^{+\infty}\frac{x^{q-1}}{(1+x)^{p+q}}\mathrm dx\quad(p,q>0).
\end{gathered}
\]

\noindent\textbf{证明}\quad (1) 令 $x=\sin^2\theta$ 作积分变换即可。

(2) 令 $x=\frac{t}{1+t}$ 作积分变换即可. 证毕。

\subsection{$\Gamma$ 函数与 B 函数的关系}

\noindent\textbf{定理 17.3.7}\quad 设 $p>0,q>0$, 则有
\[
  B(p,q)=\frac{\Gamma(p)\Gamma(q)}{\Gamma(p+q)}.
\]

\noindent\textbf{证明}\quad 由性质 17.3.2 有
\[
  \Gamma(p)=2\int_0^{+\infty}x^{2p-1}e^{-x^2}\mathrm dx,\quad
  \Gamma(q)=2\int_0^{+\infty}y^{2q-1}e^{-y^2}\mathrm dy.
\]
若令 $D=\{(x,y):0\le x<+\infty,0\le y<+\infty\}$, 则有
\[
  \Gamma(p)\Gamma(q)
  =4\iint_D x^{2p-1}y^{2q-1}e^{-(x^2+y^2)}\mathrm dx\mathrm dy.
\]
利用极坐标变换 $x=r\cos\theta,\ y=r\sin\theta$, 并令
\[
  D_1=\{(r,\theta):0<r<+\infty,0\le\theta\le\frac\pi2\},
\]
则
\[
\begin{aligned}
  \Gamma(p)\Gamma(q)
  &=4\iint_{D_1}r^{2(p+q)-1}e^{-r^2}
    \cos^{2p-1}\theta\sin^{2q-1}\theta\,\mathrm dr\mathrm d\theta\\
  &=
  \left(2\int_0^{\frac\pi2}\cos^{2p-1}\theta\sin^{2q-1}\theta\,\mathrm d\theta\right)
  \left(2\int_0^{+\infty}r^{2(p+q)-1}e^{-r^2}\mathrm dr\right)\\
  &=B(p,q)\Gamma(p+q).
\end{aligned}
\]
证毕。

\noindent\textbf{推论}\quad $B(p,q)$ 在 $(0,+\infty)\times(0,+\infty)$ 内具有任意阶偏导数。

\noindent\textbf{定理 17.3.8（余元公式）}\quad 设 $0<p<1$, 则有
\[
  B(p,1-p)=\Gamma(p)\Gamma(1-p)=\frac{\pi}{\sin p\pi}.
\]

\noindent\textbf{证明}\quad 由于
\[
  B(p,1-p)=\int_0^{+\infty}\frac{x^{p-1}}{1+x}\mathrm dx,
\]
利用变量替换 $x=\frac1t$ 有
\[
  \int_1^{+\infty}\frac{x^{p-1}}{1+x}\mathrm dx
  =\int_0^1\frac{x^{-p}}{1+x}\mathrm dx.
\]
因此
\[
  B(p,1-p)=\int_0^1\frac{x^{p-1}+x^{-p}}{1+x}\mathrm dx.
\]
我们将 $\frac1{1+x}$ 展成幂级数, 从而有
\[
\begin{aligned}
  B(p,1-p)
  &=\lim_{r\to1-0}\int_0^r\frac{x^{p-1}+x^{-p}}{1+x}\mathrm dx\\
  &=\lim_{r\to1-0}\int_0^r
  \left[\sum_{k=0}^{+\infty}(-1)^k x^{k+p-1}
  +\sum_{k=0}^{+\infty}(-1)^k x^{k-p}\right]\mathrm dx\\
  &=\sum_{k=0}^{+\infty}\frac{(-1)^k}{k+p}
  +\sum_{k=0}^{+\infty}\frac{(-1)^k}{k-p+1}\\
  &=\frac1p+\sum_{k=1}^{+\infty}(-1)^k
  \left(\frac1{k+p}+\frac1{p-k}\right)\\
  &=\frac1p+\sum_{k=1}^{+\infty}(-1)^k\frac{2p}{p^2-k^2}.
\end{aligned}
\]
由于 $\cos px$ 的傅里叶级数
\[
  \cos px=\frac{\sin p\pi}{\pi}
  \left[\frac1p+\sum_{k=1}^{+\infty}(-1)^k\frac{2p}{p^2-k^2}\cos kx\right]
\]
在 $|x|\le\pi$ 处处收敛, 令 $x=0$ 即得
\[
  B(p,1-p)=\frac1p+\sum_{k=1}^{+\infty}(-1)^k\frac{2p}{p^2-k^2}
  =\frac{\pi}{\sin p\pi}.
\]
证毕。

特别地, 在余元公式中令 $p=\frac12$, 得
\[
  B\left(\frac12,\frac12\right)=2\int_0^{\frac\pi2}\mathrm d\theta=\pi.
\]
又
\[
  \Gamma^2\left(\frac12\right)
  =\frac{\Gamma\left(\frac12\right)\Gamma\left(\frac12\right)}
  {\Gamma\left(\frac12+\frac12\right)}
  =B\left(\frac12,\frac12\right)=\pi\quad(\Gamma(1)=1),
\]
即
\[
  \Gamma\left(\frac12\right)=\sqrt\pi.
\]
由递推公式及余元公式, 若我们知道 $\Gamma(s)$ 在 $\left(0,\frac12\right)$ 中的值, 则可求出它在 $(0,+\infty)$ 的所有值. 值得注意的是, 我们又一次得到了
\[
  \int_0^{+\infty}e^{-x^2}\mathrm dx=\frac12\Gamma\left(\frac12\right)=\frac{\sqrt\pi}{2}.
\]

\noindent\textbf{例 17.3.1}\quad 计算定积分
\[
  I_n=\int_0^{\frac\pi2}\sin^n x\,\mathrm dx
  =\int_0^{\frac\pi2}\cos^n x\,\mathrm dx.
\]

\noindent\textbf{解}\quad 由 B 函数的定义我们有
\[
\begin{aligned}
  I_n
  &=\frac12\cdot2\int_0^{\frac\pi2}\cos^{\frac12\cdot2-1}x
    \sin^{\frac{2(n+1)}2-1}x\,\mathrm dx\\
  &=\frac12 B\left(\frac12,\frac{n+1}{2}\right)
  =\frac{\Gamma\left(\frac12\right)\Gamma\left(\frac{n+1}{2}\right)}
  {2\Gamma\left(\frac12+\frac{n+1}{2}\right)}.
\end{aligned}
\]
因此
\[
I_n=
\begin{cases}
  \displaystyle \frac{(n-1)!!}{n!!}\cdot\frac\pi2, & n\text{ 为偶数},\\[1.2ex]
  \displaystyle \frac{(n-1)!!}{n!!}, & n\text{ 为奇数}.
\end{cases}
\]

\noindent\textbf{例 17.3.2}\quad 求 $\mathbb R^n$ 中单位球体 $D:x_1^2+x_2^2+\cdots+x_n^2\le1$ 的体积。

\noindent\textbf{解}\quad 由 $n$ 重积分的几何意义知, 所求体积为
\[
  V=\idotsint_D \mathrm dx_1\mathrm dx_2\cdots\mathrm dx_n.
\]
类似于三维情形的球坐标变换, 我们令变换为
\[
\left\{
\begin{aligned}
  x_1&=r\cos\theta_1,\\
  x_2&=r\sin\theta_1\cos\theta_2,\\
  x_3&=r\sin\theta_1\sin\theta_2\cos\theta_3,\\
  &\cdots\cdots\cdots\\
  x_{n-1}&=r\sin\theta_1\sin\theta_2\cdots\sin\theta_{n-2}\cos\theta_{n-1},\\
  x_n&=r\sin\theta_1\sin\theta_2\cdots\sin\theta_{n-2}\sin\theta_{n-1},
\end{aligned}
\right.
\]
其中 $0<r<1,\ 0<\theta_1,\theta_2,\cdots,\theta_{n-2}<\pi,\ 0<\theta_{n-1}<2\pi$, 则其雅可比行列式
\[
  \frac{\partial(x_1,x_2,\cdots,x_n)}
  {\partial(r,\theta_1,\theta_2,\cdots,\theta_{n-2})}
  =r^{n-1}\sin^{n-2}\theta_1\sin^{n-3}\theta_2\cdots\sin\theta_{n-2}.
\]
由此得
\[
\begin{aligned}
  V
  &=\idotsint_D \mathrm dx_1\mathrm dx_2\cdots\mathrm dx_n\\
  &=\int_0^{2\pi}\mathrm d\theta_{n-1}\int_0^\pi\mathrm d\theta_{n-2}
    \cdots\int_0^\pi\mathrm d\theta_1\int_0^1
    r^{n-1}\sin^{n-2}\theta_1\sin^{n-3}\theta_2\cdots
    \sin\theta_{n-2}\mathrm dr\\
  &=\frac{2\pi}{n}
    \left(\int_0^\pi\sin^{n-2}\theta_1\mathrm d\theta_1\right)
    \left(\int_0^\pi\sin^{n-3}\theta_2\mathrm d\theta_2\right)\cdots
    \left(\int_0^\pi\sin\theta_{n-2}\mathrm d\theta_{n-2}\right)\\
  &=\frac{2\pi}{n}
    B\left(\frac12,\frac{n-1}{2}\right)
    B\left(\frac12,\frac{n-2}{2}\right)\cdots B\left(\frac12,1\right)\\
  &=\frac{2\pi}{n}\cdot
    \frac{\Gamma\left(\frac12\right)\Gamma\left(\frac{n-1}{2}\right)}{\Gamma\left(\frac n2\right)}
    \cdot
    \frac{\Gamma\left(\frac12\right)\Gamma\left(\frac{n-2}{2}\right)}{\Gamma\left(\frac{n-1}{2}\right)}
    \cdots
    \frac{\Gamma\left(\frac12\right)\Gamma(1)}{\Gamma\left(\frac32\right)}\\
  &=\frac{2\pi}{n}\cdot
    \frac{\Gamma^{n-2}\left(\frac12\right)}{\Gamma\left(\frac n2\right)}
  =\frac{\pi^{\frac n2}}{\Gamma\left(\frac n2+1\right)}.
\end{aligned}
\]

\section*{习题十七}

1. 举例说明在 $D=[0,1]\times[0,1]$ 上存在函数 $f(x,y)$, 使得它同时满足以下条件:

(1) $f(x,y)$ 的不连续点在 $D$ 稠密;

(2) 对于 $\forall x\in[0,1]$, $I(x)=\int_0^1f(x,y)\mathrm dy$ 存在, 并且 $I(x)$ 在 $[0,1]$ 上连续。

2. 设函数 $f(x,y)$ 在 $D=[a,b]\times[c,d]$ 上连续, 证明: 对于 $\forall(x,u)\in D$, $I(x,u)=\int_c^u f(x,y)\mathrm dy$ 存在, 并且 $I(x,u)$ 在 $D$ 上连续。

3. 设 $N(0,1)\subset\mathbb R^n$, 函数 $f(x,y)$ 在 $\overline{N(0,1)}\times\overline{N(0,1)}$ 上连续, 证明: 对于 $\forall x\in\overline{N(0,1)}$,
\[
  I(x)=\idotsint_{N(0,1)}f(x,y)\mathrm dy_1\mathrm dy_2\cdots\mathrm dy_n
\]
存在, 并且在 $\overline{N(0,1)}$ 上连续。

4. 求下列极限:
\[
  (1)\quad \lim_{x\to0}\int_0^{e^x}\frac{\cos xy}{\sqrt{x^2+y^2+1}}\mathrm dy;
  \qquad
  (2)\quad \lim_{x\to1}\int_0^1\frac{\mathrm dy}{1+xy^2}.
\]

5. 计算无穷积分
\[
  \int_0^{+\infty}\frac{e^{-ax}-e^{-bx}}{x}\mathrm dx\quad(a>b>0).
\]

6. 计算无穷积分
\[
  \int_0^{+\infty}\frac{e^{-ax^2}-e^{-bx^2}}{x^2}\mathrm dx\quad(b>a>0).
\]

\end{document}
